\documentclass{article}
\usepackage{mathtools} 
\usepackage{fontspec}
\usepackage[UTF8]{ctex}
\usepackage{amsthm}
\usepackage{mdframed}
\usepackage{xcolor}
\usepackage{amssymb}
\usepackage{amsmath}


% 定义新的带灰色背景的说明环境 zremark
\newmdtheoremenv[
  backgroundcolor=gray!10,
  % 边框与背景一致，边框线会消失
  linecolor=gray!10
]{zremark}{说明}

% 通用矩阵命令: \flexmatrix{矩阵名}{元素符号}{行数}{列数}
\newcommand{\flexmatrix}[4]{
  \[
  #1 = \begin{pmatrix}
    #2_{11}     & #2_{12}     & \cdots & #2_{1#4}   \\
    #2_{21}     & #2_{22}     & \cdots & #2_{2#4}   \\
    \vdots      & \vdots      & \ddots & \vdots     \\
    #2_{#31}    & #2_{#32}    & \cdots & #2_{#3#4}
  \end{pmatrix}
  \]
}

% 简化版命令（默认矩阵名为A，元素符号为a）: \quickmatrix{行数}{列数}
\newcommand{\quickmatrix}[2]{\flexmatrix{A}{a}{#1}{#2}}

\begin{document}
\title{总结：第二章}
\author{张志聪}
\maketitle

\section*{1}

\begin{zremark}
  核心：

  1.计算极限的方式。
\end{zremark}

\begin{itemize}
  \item 通过定义$\epsilon - N$方式；
  \item 迫敛性；
  \item 四则运算；
  \item 子列；
  \item 单调有界定理；
  \item 柯西收敛准则；
\end{itemize}

\begin{zremark}
  几个重要数列极限
\end{zremark}

\begin{itemize}
  \item $\lim\limits_{n \to \infty} \frac{1}{n^\alpha} = 0$，其中$\alpha > 0$（$\S 1$例2）；
  \item $\lim\limits_{n \to \infty} q^n = 0$，其中$|q| < 1$（$\S 1$例4）；
  \item $\lim\limits_{n \to \infty} \sqrt[n]{a} = 1$，其中$a > 0$（$\S 1$例5）；
  \item $\lim\limits_{n \to \infty} \frac{a^n}{n!} = 0$（$\S 1$例6）；
  \item $\lim\limits_{n \to \infty} \frac{n!}{n^n} = 0$（$\S 1$习题2-3）；
  \item $\lim\limits_{n \to \infty} \frac{n}{a^n} = 0$（$\S 1$习题2-5）；
  \item $\lim\limits_{n \to \infty} \sqrt[n]{n} = 1$（$\S 2$例2）；
  \item $\lim\limits_{n \to \infty} \frac{1}{\sqrt[n]{n!}} = 0$（$\S 2$例3）；
\end{itemize}

\begin{zremark}
  对数列的项进行处理的常用方式。
\end{zremark}

\begin{itemize}
  \item 分式，分母缩小。（$\S 1$例3）（\S2 例1）。
  \item $n$次幂降幂（$\S 1$例4），运用$(1+h)^n \geq 1 + nh$
        或使用二项展开$(1+h)^n \geq 1 + nh + \frac{n(n-1)}{2}h^2$（\S2 例2），进行降幂。
  \item $\frac{1}{n}$次幂，升幂（$\S 1$例5）。

        先升到1次幂：

        令
        \begin{align*}
          \alpha_n = a^\frac{1}{n} - a
        \end{align*}

        然后
        \begin{align*}
          \alpha_n + a = a^\frac{1}{n} \\
          (\alpha_n + a)^n = a
        \end{align*}

        左侧降幂

  \item 累积形式（$\S 1$ 例6），以累积中的合适项，代替其他项或部分项，进行放大。
  \item 累积形式（$\S 2$ 习题8-1），插入项，进行放大或缩小。
  \item 累加（或累乘）形式表示$a_n$，用特定项（最小项或最大项或某一项）代替所有项（或部分项），进行缩放。
        如：$\S 2$习题4-(5)，$\S 2$习题8-(2)
  \item 累加形式，按奇数偶数拆开。
  \item 累加形式，分母裂项。如：$\S 3$习题5-(2)
  \item 累乘利用$n^2 \geq n^2 - 1 = (n - 1)(n + 1)$不等式，对分母进行拆分，达到与分子约分的目的。
        如：$\S 2$习题8-(1)
  \item 加项，使得$a_n$可以化简。如：$\S 2$习题8-(4)
  \item 利用已有的极限（$\S2$例3）。
\end{itemize}

\begin{zremark}
  利用“单调有界定理”时，如何证明数列$\{a_n\}$单调性和有界性。
\end{zremark}

单调性
\begin{itemize}
  \item 相减
  \item 相除
  \item 原始的$a_n, a_{n + 1}$比较起来相对困难，修改$a_n, a_{n + 1}$的表示方式，便于比较（$\S 3$例4）。
\end{itemize}

有界性
\begin{itemize}
  \item 由单调性（以单增为例）可知，$a_n \leq a_{2n}$，然后用$a_n$表示出$a_{2n} < a + k a_{n} \ (k < 1)$，于是可得
        $a_n < a + k a_n$得出$a_n$的范围。（$\S 3$例1）
\end{itemize}

\begin{zremark}
  利用已知的数列极限的方式：
  $\left(1 + \frac{1}{n^2}\right)^{n^2}$是
  $\left(1 + \frac{1}{n}\right)^{n} = e $的子列
  可得
  \begin{align*}
    \left(1 + \frac{1}{n^2}\right)^{n^2} = e
  \end{align*}
\end{zremark}

\end{document}